\documentclass[11pt]{article}
\usepackage[margin=1in]{geometry}
\usepackage{amsmath}

\title{A Tiny Example for \texttt{pdf-review}}
\author{}
\date{}

\begin{document}
\maketitle

\begin{abstract}
A small document for exercising \texttt{pdf-review}. It has prose here in
\texttt{main.tex} and four sections pulled in from separate files, together with
equations, tables and enough text to run over several pages, so that clicking
around the compiled PDF shows whether a selection is resolved to the right source
file and not merely to the right line. The content is deliberately unremarkable:
a constant-returns production function, a household problem, a short panel of
invented data, and the regressions one would run on it. Nothing in the argument
needs defending, which means any passage can be marked, rewritten or deleted
without the test losing its point. What the document is arranged to stress is
everything between a rendered box and the line that produced it: paragraphs that
wrap across a page break, displayed equations that sit a line away from the text
introducing them, floating tables that \LaTeX{} is free to move far from their
source, and \verb|\paragraph| headings that put several plausible lines within a
sentence of each other. A selection that survives all four, in the right file,
is a selection an editing agent can act on.
\end{abstract}

\section{Introduction}

This document exists so you can try \texttt{pdf-review} end to end. Compile it
with SyncTeX enabled, start the viewer, and Cmd/Ctrl-click anywhere below: the
source file and line for that spot are written to
\texttt{.pdf-review/current-selection.json}.

The interesting case is text that lives in another file, so the model section is
pulled in with \verb|\input|. The same goes for the data, results and robustness
sections: each one is its own file under \texttt{sections/}, and each one is a
different answer to the question ``where did that sentence come from?''

To make the exercise genuinely useful, the document is long enough to spill over
several pages. Page breaks are where a naive line-based mapping tends to go
wrong, because the box a reader clicks on may have been set on one page while
the paragraph that produced it started on the previous one. Floating tables make
this worse, since \LaTeX{} is free to move them far from the source line that
defined them.

What follows is vanilla economics by design: a production function, a
household problem, a panel of invented data, and the regressions one would run
on it. Ordinary is the point. A reviewer's attention should go to the machinery
for marking and rewriting passages, not to the argument, and any sentence here
can be replaced without anyone mourning it.

\section{Model}

Let output be produced by a constant-returns technology, so that
\begin{equation}
  Y_{t} = A_{t} K_{t}^{\alpha} L_{t}^{1 - \alpha},
  \label{eq:production}
\end{equation}
where $A_{t}$ is total factor productivity, $K_{t}$ is capital, and $L_{t}$ is
labour. The elasticity $\alpha$ is the parameter a reader is most likely to
argue with, which makes this a good paragraph to click on.

Equation~\eqref{eq:production} is deliberately unremarkable. What matters is
that this sentence lives in \texttt{sections/model.tex} rather than in
\texttt{main.tex}, so a click here exercises the path mapping back to the
project root.

\subsection{Households}

A representative household chooses consumption $C_{t}$ and next-period capital
$K_{t+1}$ to maximise
\begin{equation}
  \sum_{t = 0}^{\infty} \beta^{t} \, u(C_{t}),
  \qquad
  u(C) = \frac{C^{1 - \sigma} - 1}{1 - \sigma},
  \label{eq:utility}
\end{equation}
subject to the resource constraint
\begin{equation}
  C_{t} + K_{t+1} = Y_{t} + (1 - \delta) K_{t},
  \label{eq:budget}
\end{equation}
with discount factor $\beta \in (0, 1)$, intertemporal elasticity $1 / \sigma$,
and depreciation rate $\delta$. Labour is supplied inelastically at $L_{t} = 1$,
which is the sort of assumption that survives in a draft precisely because
nobody bothers to click on it.

Combining the first-order conditions for~\eqref{eq:utility} and~\eqref{eq:budget}
means using the budget constraint to write
$C_{t} = Y_{t} + (1 - \delta) K_{t} - K_{t+1}$, substituting it into the
objective, and differentiating with respect to the single remaining choice
variable $K_{t+1}$. Saving one more unit of capital costs one unit of consumption
today, which is worth $\beta^{t} u'(C_{t})$. Tomorrow that unit pays back two
things: the extra output it produces, $\alpha A_{t+1} K_{t+1}^{\alpha - 1}$,
which is the marginal product implied by~\eqref{eq:production}, and the
$1 - \delta$ units of capital left over after depreciation. Its gross return is
therefore $\alpha A_{t+1} K_{t+1}^{\alpha - 1} + 1 - \delta$, and the payoff is
worth $\beta^{t+1} u'(C_{t+1})$ times that return. Setting the cost equal to the
payoff and using $u'(C) = C^{-\sigma}$ gives the Euler equation
\begin{equation}
  \left( \frac{C_{t+1}}{C_{t}} \right)^{\sigma}
    = \beta \left[ \alpha A_{t+1} K_{t+1}^{\alpha - 1} + 1 - \delta \right],
  \label{eq:euler}
\end{equation}
which says that the utility the household gives up by saving one more unit
today must exactly match the discounted utility that unit returns tomorrow.
Everything that follows is a consequence of~\eqref{eq:euler} and a
transversality condition we will not write down.

\subsection{Steady state}

Setting $A_{t} = A$ and $C_{t+1} = C_{t}$ in~\eqref{eq:euler} pins down the
capital stock,
\begin{equation}
  K^{\star}
    = \left( \frac{\alpha A}{\beta^{-1} - 1 + \delta} \right)^{\frac{1}{1 - \alpha}},
  \label{eq:kstar}
\end{equation}
and steady-state consumption follows from~\eqref{eq:budget} as
$C^{\star} = A (K^{\star})^{\alpha} - \delta K^{\star}$. The comparative statics
are the ones every reader expects: $K^{\star}$ rises with patience and with
productivity, and falls as capital wears out faster.

Two features of~\eqref{eq:kstar} matter later. First, the exponent
$1 / (1 - \alpha)$ means that a small change in $\alpha$ moves the steady state a
lot, so the calibration of $\alpha$ deserves more attention than it usually
receives. Second, nothing in the steady state depends on $\sigma$; the
intertemporal elasticity governs the speed of transition rather than the
destination. Readers who dislike this separation tend to say so in the margin,
which is exactly the kind of comment this document is built to collect.


\section{Robustness}
\label{sec:robustness}

Three assumptions carry more weight than the rest, and each is worth a paragraph
that a reader can disagree with directly.

\paragraph{Depreciation.} Setting $\delta = 0.06$ is conventional rather than
estimated. Re-running the perpetual inventory recursion~\eqref{eq:pim} at
$\delta = 0.04$ and $\delta = 0.08$ shifts the level of the capital series
noticeably but leaves $\Delta \log K$ nearly untouched, so $\hat{\alpha}$ moves
by less than a point. The capital-output ratio in Table~\ref{tab:summary} is not
so lucky: it ranges from $2.4$ to $3.4$ across the three choices, and since that
moment is what identifies $\alpha$ in the calibrated version of the model, the
two routes to $\alpha$ disagree by more than either one's standard error admits.

\paragraph{Constant returns.} Imposing $\alpha + (1 - \alpha) = 1$
in~\eqref{eq:production} is what allows the factor shares in
Table~\ref{tab:estimates} to be read off a single coefficient. Relaxing it and
estimating the two elasticities freely gives $0.36$ and $0.61$, whose sum is
statistically indistinguishable from one and economically indistinguishable from
anything else in the neighbourhood. This is reassuring without being evidence.

\subsection{Measurement of the residual}

Because $\Delta \log A_{t}$ is defined as what~\eqref{eq:growth} does not
explain, every measurement error in output, capital or employment lands in it
with the wrong sign attached. The persistence of the estimated residual is
therefore an upper bound on the persistence of true productivity, and the
variance decomposition in Section~\ref{sec:data} should be read with that in
mind.

None of this changes the qualitative picture, which is the sentence every
robustness section ends on and the one readers most enjoy clicking to complain
about. What it does change is how much confidence the twenty per cent band
around $K^{\star}$ deserves, and the honest answer is: less than the tables
imply, more than a sceptic would grant.


\section{Data}
\label{sec:data}

The exercise uses an unbalanced panel of $N = 42$ imaginary economies observed
annually from 1980 to 2019. Output, capital and employment are the three series
we need; everything else in the source files is carried along because deleting
columns from a panel is the kind of chore that always waits until the referee
report arrives.

Capital is constructed by the perpetual inventory method,
\begin{equation}
  K_{t+1} = (1 - \delta) K_{t} + I_{t},
  \label{eq:pim}
\end{equation}
initialised at the average investment-output ratio of the first five years and a
depreciation rate of $\delta = 0.06$. The initial condition matters much less
than one might fear: by construction the influence of $K_{0}$ decays at rate
$(1 - \delta)^{t}$, so after two decades roughly three per cent of it survives.

\begin{table}[ht]
  \centering
  \begin{tabular}{lrrrr}
    \hline
    Variable & Mean & Std.\ dev. & Min & Max \\
    \hline
    Output growth, $\Delta \log Y$   & 0.024 & 0.031 & $-0.084$ & 0.112 \\
    Capital growth, $\Delta \log K$  & 0.031 & 0.019 & $-0.012$ & 0.094 \\
    Employment growth, $\Delta \log L$ & 0.009 & 0.014 & $-0.041$ & 0.058 \\
    Investment rate, $I / Y$         & 0.221 & 0.047 & 0.108 & 0.361 \\
    Capital-output ratio, $K / Y$    & 2.870 & 0.612 & 1.640 & 4.410 \\
    \hline
  \end{tabular}
  \caption{Summary statistics for the estimation sample, 1980--2019.}
  \label{tab:summary}
\end{table}

Table~\ref{tab:summary} reports the moments the model is asked to match. The
capital-output ratio is the one that does the most work, since together
with~\eqref{eq:kstar} it identifies $\alpha$ up to the assumed real interest
rate. Nothing here is real, but the numbers are chosen to look plausible enough
that a reader will click on them before checking whether they are.

Two sample restrictions are worth flagging. Economies with fewer than fifteen
consecutive observations are dropped, which removes six of the smallest units
and about four per cent of the observations. Years in which measured output
falls by more than ten per cent are kept rather than winsorised; they are rare,
they are informative about the residual, and dropping them would be the sort of
quiet decision that a reader has every right to object to in the margin.


\section{Results}

Taking logs of~\eqref{eq:production} and differencing over time gives the usual
growth-accounting decomposition,
\begin{equation}
  \Delta \log Y_{t} = \Delta \log A_{t} + \alpha \, \Delta \log K_{t}
    + (1 - \alpha) \, \Delta \log L_{t},
  \label{eq:growth}
\end{equation}
so measured output growth splits into a productivity term and the two factor
terms weighted by $\alpha$ and $1 - \alpha$.

The result worth stating is what~\eqref{eq:growth} leaves behind: whatever is not
explained by capital and labour is charged to $A_{t}$, which is why the residual
carries so much of the action in practice. This paragraph also lives in
\texttt{sections/results.tex}, giving a second file to click on.

\subsection{Point estimates}

Estimating~\eqref{eq:growth} on the panel of Section~\ref{sec:data} by
least squares with unit and year effects gives $\hat{\alpha} = 0.34$ with a
standard error of $0.02$, comfortably inside the range the calibration
literature treats as uncontroversial. The implied productivity residual accounts
for a little under half of average output growth, and rather more than half of
its variance across economies.

\begin{table}[ht]
  \centering
  \begin{tabular}{lrrr}
    \hline
    & \multicolumn{1}{c}{(1)} & \multicolumn{1}{c}{(2)} & \multicolumn{1}{c}{(3)} \\
    \hline
    $\Delta \log K$  & 0.34 & 0.31 & 0.36 \\
                     & (0.02) & (0.03) & (0.02) \\
    $\Delta \log L$  & 0.66 & 0.69 & 0.64 \\
                     & (0.02) & (0.03) & (0.02) \\
    Unit effects     & Yes & Yes & Yes \\
    Year effects     & No  & Yes & Yes \\
    Balanced sample  & No  & No  & Yes \\
    \hline
    Observations     & 1{,}512 & 1{,}512 & 1{,}280 \\
    $R^{2}$          & 0.61 & 0.68 & 0.71 \\
    \hline
  \end{tabular}
  \caption{Growth-accounting regressions. Standard errors clustered by unit in
    parentheses; factor shares constrained to sum to one.}
  \label{tab:estimates}
\end{table}

Table~\ref{tab:estimates} shows that the estimate is not delicate. Adding year
effects moves $\hat{\alpha}$ by three points, restricting attention to the
balanced sample moves it by two in the other direction, and no specification
takes it outside $[0.30, 0.37]$. Given the exponent in~\eqref{eq:kstar}, however,
even that narrow band maps into a range for $K^{\star}$ of roughly twenty per
cent, which is a good deal less comforting than the standard errors suggest.

\subsection{Transition dynamics}

Simulating the model forward from a capital stock half its steady-state value
produces the familiar concave path: output growth starts near six per cent and
decays geometrically toward zero. The half-life of the gap is
\begin{equation}
  h = \frac{\log 2}{\lambda},
  \qquad
  \lambda \approx (1 - \alpha)\left( \beta^{-1} - 1 + \delta \right),
  \label{eq:halflife}
\end{equation}
which at the estimated parameters is a little over eighteen years. Convergence
in this class of models is slow, and~\eqref{eq:halflife} makes clear that it is
slow for a structural reason rather than an accident of calibration: raising
$\alpha$ lowers $\lambda$ directly.

The gap between the model's transition and the data's is where an honest draft
starts hedging. The simulated paths are smoother than the observed ones, they
lack the persistence of the measured residual, and they cannot reproduce the
handful of large contractions that survive in the sample. A reader who wants to
argue with any of that will find this paragraph a convenient place to start.


\section{Conclusion}

Poke this paragraph and it should shout \texttt{main.tex}; poke the model section
and it should shout \texttt{sections/model.tex}. That is the whole trick: an
agent handed a naked line number is a tourist without a map.

With five sections spread across four files, the mapping has room to fail in the
ways that matter. A click in the middle of a wrapped paragraph should still land
on the paragraph, a click in a displayed equation should land on the equation
rather than on the line after it, and a click inside a table should land
somewhere in the \texttt{tabular} rather than wherever the float happened to
drift to.

Section~\ref{sec:robustness} is the best place to try the awkward cases, since
its \verb|\paragraph| headings put short runs of text next to each other and
give SyncTeX several plausible lines to choose from. If a selection there comes
back with the right file and a line within a sentence or two of the text you
meant, the loop is working as intended; the model reads the surrounding lines
anyway and finds the passage from the quoted words.

\end{document}
